\chapter{Capitolo 12 - Protocollo scientifico OCT: metrica, osservazione, decisione}

\section{1. Problema}

Con i capitoli precedenti abbiamo costruito un impianto teorico completo:

\begin{enumerate}
\def\labelenumi{\arabic{enumi}.}
\tightlist
\item
  fondazione ontologica e assiomatica;
\item
  teoremi fondativi, differenziali e applicativi;
\item
  estensione del corpus classico fino alle strutture avanzate;
\item
  grammatica di controesempi, limiti e falsificabilità.
\end{enumerate}

Manca però il passaggio che decide se la teoria può diventare pratica scientifica condivisa: un protocollo unico che dica \textbf{come} osservare, \textbf{cosa} misurare e \textbf{quando} decidere.

Senza protocollo:

\begin{enumerate}
\def\labelenumi{\arabic{enumi}.}
\tightlist
\item
  i risultati restano locali e difficili da confrontare;
\item
  le decisioni di validazione diventano opache;
\item
  i passaggi \texttt{in\_proof\ -\textgreater{}\ validated} rischiano arbitrarietà.
\end{enumerate}

Con protocollo:

\begin{enumerate}
\def\labelenumi{\arabic{enumi}.}
\tightlist
\item
  ogni test diventa replicabile;
\item
  ogni decisione è tracciabile;
\item
  la teoria può essere auditata da team indipendenti.
\end{enumerate}

Il problema del capitolo è quindi:

\textbf{come formalizzare una procedura standard che unisca metrica ordinativa, osservazione contestuale e decisione epistemica senza ridurre OCT a sola statistica né a sola deduzione formale?}

\subsection{\texorpdfstring{1.1 Perché questo capitolo apre la milestone \texttt{v1.4}}{1.1 Perché questo capitolo apre la milestone v1.4}}

La milestone \texttt{v1.4} è il passaggio da ``teoria estesa'' a ``metodo scientifico operativo'':

\begin{enumerate}
\def\labelenumi{\arabic{enumi}.}
\tightlist
\item
  Capitolo 12: protocollo;
\item
  Capitolo 13: benchmark e replicabilità;
\item
  Capitolo 14: governance epistemica.
\end{enumerate}

Il Capitolo 12 è la base perché definisce il linguaggio procedurale comune usato nei due capitoli successivi.

\subsection{1.2 Vincolo metodologico}

Il protocollo OCT deve mantenere una tripla coerenza:

\begin{enumerate}
\def\labelenumi{\arabic{enumi}.}
\tightlist
\item
  coerenza matematica (con O1-O7 e i teoremi F/D/A);
\item
  coerenza sperimentale (misure replicabili);
\item
  coerenza decisionale (regole esplicite di accettazione/revisione/rigetto).
\end{enumerate}

\begin{center}\rule{0.5\linewidth}{0.5pt}\end{center}

\section{2. Tesi del capitolo}

La tesi è articolata in nove enunciati.

\begin{enumerate}
\def\labelenumi{\arabic{enumi}.}
\tightlist
\item
  Un risultato OCT è scientificamente forte solo se combina prova formale e protocollo di osservazione replicabile.
\item
  La triade metrica canonica (\texttt{Coh}, \texttt{Phi}, \texttt{Delta}) è necessaria ma non sufficiente: servono metriche ausiliarie dominio-specifiche.
\item
  Ogni esperimento deve essere indicizzato da un contesto \texttt{Omega} dichiarato e versionato.
\item
  La pre-registrazione dei criteri decisionali riduce bias di conferma.
\item
  La validazione richiede benchmark indipendenti e multi-contesto.
\item
  I criteri di decisione devono distinguere \texttt{validated}, \texttt{revise}, \texttt{reject}.
\item
  Il protocollo deve prevedere pubblicazione di fallimenti significativi, non solo esiti positivi.
\item
  Il passaggio di stato di un teorema è legittimo solo con evidenza tracciata.
\item
  Il protocollo proposto è sufficiente per sostenere il disegno benchmark del Capitolo 13.
\end{enumerate}

\begin{center}\rule{0.5\linewidth}{0.5pt}\end{center}

\section{3. Definizioni canoniche}

\subsection{Definizione 12.1 (Unità di validazione)}

Unità minima di validazione:

\texttt{U\_val\ :=\ \textless{}Claim,\ Omega,\ Dataset,\ Pipeline,\ Metriche,\ Criteri,\ Log,\ Esito\textgreater{}}.

\subsection{Definizione 12.2 (Contesto osservativo registrato)}

\texttt{Omega\_reg} è il contesto osservativo con:

\begin{enumerate}
\def\labelenumi{\arabic{enumi}.}
\tightlist
\item
  variabili rilevanti;
\item
  ipotesi di dominio;
\item
  condizioni di perturbazione ammesse;
\item
  versione del protocollo.
\end{enumerate}

\subsection{Definizione 12.3 (Metriche canoniche e ausiliarie)}

Metriche canoniche:

\begin{enumerate}
\def\labelenumi{\arabic{enumi}.}
\tightlist
\item
  \texttt{Coh\_Omega(D)\ in\ {[}0,1{]}};
\item
  \texttt{Phi\_Omega(D)} (normalizzata quando necessario);
\item
  \texttt{Delta\_Omega(D)=1-Coh\_Omega(D)}.
\end{enumerate}

Metriche ausiliarie (esempi):

\begin{enumerate}
\def\labelenumi{\arabic{enumi}.}
\tightlist
\item
  \texttt{Err\_sem} per qualità semantica;
\item
  \texttt{Err\_struct} per fedeltà strutturale;
\item
  \texttt{Var\_Omega} per stabilità cross-contesto.
\end{enumerate}

\subsection{Definizione 12.4 (Pre-registrazione)}

La pre-registrazione è la dichiarazione ex ante di:

\begin{enumerate}
\def\labelenumi{\arabic{enumi}.}
\tightlist
\item
  criteri di conferma;
\item
  criteri di falsificazione;
\item
  soglie e test statistici;
\item
  regole di decisione.
\end{enumerate}

\subsection{Definizione 12.5 (Decisione epistemica)}

Esito finale di un test:

\begin{enumerate}
\def\labelenumi{\arabic{enumi}.}
\tightlist
\item
  \texttt{validated}: evidenza sufficiente e replicata;
\item
  \texttt{revise}: evidenza mista o insufficiente;
\item
  \texttt{reject}: controevidenza robusta nel dominio dichiarato.
\end{enumerate}

\subsection{Definizione 12.6 (Tracciabilità completa)}

Un risultato è tracciabile se include:

\begin{enumerate}
\def\labelenumi{\arabic{enumi}.}
\tightlist
\item
  configurazione di esecuzione;
\item
  dati/versione;
\item
  codice o procedura;
\item
  output e log;
\item
  motivazione della decisione.
\end{enumerate}

\subsection{Definizione 12.7 (Indice di evidenza composta)}

Dato un claim \texttt{C} e un contesto \texttt{Omega}, definiamo l'indice di evidenza composta:

\texttt{IEC(C,Omega)\ :=\ w1*Coh\ +\ w2*Phi\_norm\ +\ w3*(1-Delta)\ +\ w4*Rep\ +\ w5*Trace}

con:

\begin{enumerate}
\def\labelenumi{\arabic{enumi}.}
\tightlist
\item
  \texttt{w\_i\ \textgreater{}=\ 0} e \texttt{sum\_i\ w\_i\ =\ 1};
\item
  \texttt{Phi\_norm} normalizzata nel dominio del benchmark;
\item
  \texttt{Rep} indice di replicabilità inter-benchmark;
\item
  \texttt{Trace} indice di completezza di tracciabilità.
\end{enumerate}

L'\texttt{IEC} non sostituisce la valutazione teorica, ma fornisce una sintesi comparabile tra esperimenti.

\subsection{Definizione 12.8 (Budget d'incertezza)}

Per ogni unità \texttt{U\_val} definiamo un budget d'incertezza:

\texttt{B\_unc\ :=\ \textless{}u\_data,\ u\_model,\ u\_measure,\ u\_context,\ u\_decision\textgreater{}}

dove ogni termine quantifica la componente di incertezza legata a:

\begin{enumerate}
\def\labelenumi{\arabic{enumi}.}
\tightlist
\item
  qualità e copertura dei dati;
\item
  specificazione del modello/pipeline;
\item
  errore o variabilità di misura;
\item
  trasferibilità tra contesti;
\item
  discrezionalità residua in decisione.
\end{enumerate}

Un claim non può essere promosso a \texttt{validated} se \texttt{B\_unc} supera la soglia preregistrata.

\begin{center}\rule{0.5\linewidth}{0.5pt}\end{center}

\section{4. Sviluppo formale}

\subsection{4.1 Architettura del protocollo in tre strati}

Il protocollo OCT è organizzato in tre strati:

\begin{enumerate}
\def\labelenumi{\arabic{enumi}.}
\tightlist
\item
  \textbf{Strato Metrico}: cosa misurare (\texttt{Coh/Phi/Delta} + ausiliarie);
\item
  \textbf{Strato Osservativo}: in quale \texttt{Omega} e con quali perturbazioni;
\item
  \textbf{Strato Decisionale}: come passare da misura a stato teorico.
\end{enumerate}

La separazione evita confusioni:

\begin{enumerate}
\def\labelenumi{\arabic{enumi}.}
\tightlist
\item
  buone metriche con cattivo contesto;
\item
  buon contesto con regole decisionali opache;
\item
  decisioni forti su evidenza debole.
\end{enumerate}

\subsection{4.2 Pipeline minima di validazione (PV-8)}

Proponiamo una pipeline in otto passi.

\begin{enumerate}
\def\labelenumi{\arabic{enumi}.}
\tightlist
\item
  dichiarare claim e teorema target;
\item
  registrare \texttt{Omega\_reg};
\item
  fissare dataset/versione e baseline;
\item
  pre-registrare criteri di conferma/falsificazione;
\item
  eseguire esperimento e raccogliere metriche;
\item
  ripetere su benchmark indipendente;
\item
  consolidare risultati multi-contesto;
\item
  assegnare esito (\texttt{validated}/\texttt{revise}/\texttt{reject}) con motivazione.
\end{enumerate}

La PV-8 è il nucleo operativo del capitolo.

\subsection{4.3 Regole di evidenza minima}

Per promuovere un claim a \texttt{validated} proponiamo quattro condizioni minime:

\begin{enumerate}
\def\labelenumi{\arabic{enumi}.}
\tightlist
\item
  conferma su almeno due benchmark indipendenti;
\item
  replicazione su almeno due contesti \texttt{Omega} non equivalenti banalmente;
\item
  coerenza tra metrica canonica e metrica ausiliaria rilevante;
\item
  assenza di controesempio forte non risolto nel dominio dichiarato.
\end{enumerate}

Se una condizione manca:

\begin{enumerate}
\def\labelenumi{\arabic{enumi}.}
\tightlist
\item
  default su \texttt{revise} (non su \texttt{validated}).
\end{enumerate}

\subsection{4.4 Trattamento dei risultati discordanti}

È frequente ottenere risultati misti tra benchmark o contesti. In OCT non si risolve con media semplice.

Regola proposta:

\begin{enumerate}
\def\labelenumi{\arabic{enumi}.}
\tightlist
\item
  identificare la fonte di divergenza (dati, mapping, soglia, pipeline);
\item
  classificare divergenza come:

  \begin{itemize}
  \tightlist
  \item
    rumore controllabile;
  \item
    instabilità strutturale;
  \item
    controevidenza sostanziale.
  \end{itemize}
\item
  decidere esito in funzione della classe di divergenza.
\end{enumerate}

Questa regola impedisce sia ottimismo ingiustificato sia rigetto prematuro.

\subsection{4.5 Matrice decisionale standard}

Per ogni unità \texttt{U\_val}:

\begin{enumerate}
\def\labelenumi{\arabic{enumi}.}
\tightlist
\item
  \textbf{D-Valid}: criteri soddisfatti, replicazione robusta;
\item
  \textbf{D-Revise}: evidenza parziale o conflittuale;
\item
  \textbf{D-Reject}: fallimento robusto nel dominio dichiarato.
\end{enumerate}

Ogni decisione deve includere:

\begin{enumerate}
\def\labelenumi{\arabic{enumi}.}
\tightlist
\item
  claim colpito;
\item
  evidenze principali;
\item
  limiti residui;
\item
  azione successiva.
\end{enumerate}

\subsection{4.6 Collegamento con la falsificabilità (Capitolo 11)}

Il protocollo incorpora i principi K1-K7:

\begin{enumerate}
\def\labelenumi{\arabic{enumi}.}
\tightlist
\item
  riproducibilità e tracciabilità;
\item
  localizzazione del fallimento;
\item
  decisione finale non arbitraria.
\end{enumerate}

In questo modo, la falsificabilità non resta solo principio teorico: diventa procedura eseguibile.

\subsection{4.7 Controllo di qualità del protocollo}

Ogni campagna sperimentale dovrebbe produrre un report QA con:

\begin{enumerate}
\def\labelenumi{\arabic{enumi}.}
\tightlist
\item
  completezza metadati;
\item
  qualità log;
\item
  copertura benchmark;
\item
  coerenza tra decisione e criteri pre-registrati.
\end{enumerate}

Se il QA è insufficiente:

\begin{enumerate}
\def\labelenumi{\arabic{enumi}.}
\tightlist
\item
  il risultato non può passare a \texttt{validated}.
\end{enumerate}

\subsection{4.8 Versionamento del protocollo}

Il protocollo deve essere versionato (\texttt{P\_v1}, \texttt{P\_v2}, \ldots). Ogni versione deve dichiarare:

\begin{enumerate}
\def\labelenumi{\arabic{enumi}.}
\tightlist
\item
  cosa cambia nelle metriche;
\item
  cosa cambia nei criteri decisionali;
\item
  quali risultati pregressi vanno rivalutati.
\end{enumerate}

Questo previene inconsistenze storiche nei passaggi di stato.

\subsection{4.9 Protocollo e apertura dei dati}

Per quanto possibile, la validazione deve favorire:

\begin{enumerate}
\def\labelenumi{\arabic{enumi}.}
\tightlist
\item
  pubblicazione di dataset/manifest;
\item
  pubblicazione dei log principali;
\item
  pubblicazione degli errori significativi.
\end{enumerate}

La trasparenza non è accessorio etico: è condizione tecnica di replicabilità.

\subsection{4.10 Condizione di non-ridondanza del capitolo}

Il capitolo è non ridondante se mostra che:

\begin{enumerate}
\def\labelenumi{\arabic{enumi}.}
\tightlist
\item
  senza protocollo la stessa teoria può produrre esiti incompatibili;
\item
  con protocollo gli esiti diventano confrontabili e decidibili.
\end{enumerate}

Questa condizione è soddisfatta dall'integrazione tra PV-8, regole di evidenza minima e matrice decisionale.

\subsection{\texorpdfstring{4.11 Operatore decisionale ordinativo \texttt{D\_OCT}}{4.11 Operatore decisionale ordinativo D\_OCT}}

Per rendere la decisione riproducibile introduciamo un operatore esplicito:

\texttt{D\_OCT(U\_val)\ -\textgreater{}\ \{validated,\ revise,\ reject\}}.

Input minimi:

\begin{enumerate}
\def\labelenumi{\arabic{enumi}.}
\tightlist
\item
  indicatori canonici e ausiliari;
\item
  esito repliche indipendenti;
\item
  \texttt{B\_unc} e \texttt{IEC};
\item
  presenza/assenza di controesempi forti.
\end{enumerate}

Regola orientativa:

\begin{enumerate}
\def\labelenumi{\arabic{enumi}.}
\tightlist
\item
  \texttt{validated} se tutte le condizioni minime sono soddisfatte, \texttt{IEC\ \textgreater{}=\ tau\_val} e \texttt{B\_unc\ \textless{}=\ beta\_val};
\item
  \texttt{revise} se evidenza intermedia, conflittuale o incertezza oltre soglia di validazione;
\item
  \texttt{reject} se controevidenza robusta o violazioni strutturali persistenti.
\end{enumerate}

L'operatore non elimina il giudizio scientifico, ma ne impone la struttura argomentativa.

\subsection{4.12 Budget d'incertezza e analisi di sensibilità}

Ogni campagna deve includere almeno tre verifiche:

\begin{enumerate}
\def\labelenumi{\arabic{enumi}.}
\tightlist
\item
  \textbf{sensibilità dati}: variazione controllata del campione/stratificazione;
\item
  \textbf{sensibilità soglie}: stress su \texttt{tau\_val}, \texttt{beta\_val}, soglie ausiliarie;
\item
  \textbf{sensibilità contesto}: cambio di \texttt{Omega} mantenendo claim invariato.
\end{enumerate}

Se l'esito cambia in modo non spiegato rispetto a piccole perturbazioni, il claim va in \texttt{revise} fino a chiarimento.

\subsection{4.13 Schema minimo di preregistrazione}

Per evitare preregistrazioni deboli proponiamo uno schema minimo obbligatorio.

Campi obbligatori:

\begin{enumerate}
\def\labelenumi{\arabic{enumi}.}
\tightlist
\item
  claim e dominio di validità dichiarato;
\item
  dataset, fonti, criteri di inclusione/esclusione;
\item
  pipeline esecutiva e versioni software;
\item
  metriche canoniche e metriche ausiliarie;
\item
  soglie decisionali (\texttt{tau\_val}, \texttt{tau\_rev}, \texttt{beta\_val});
\item
  test di robustezza previsti;
\item
  criteri espliciti di falsificazione;
\item
  politica di gestione risultati inattesi.
\end{enumerate}

Senza questi campi, l'unita non e idonea a validazione conclusiva.

\subsection{4.14 Livelli di riproducibilità degli artefatti}

Definiamo quattro livelli progressivi:

\begin{enumerate}
\def\labelenumi{\arabic{enumi}.}
\tightlist
\item
  \texttt{R0} - descrizione narrativa non eseguibile;
\item
  \texttt{R1} - script disponibili ma ambiente non congelato;
\item
  \texttt{R2} - script + ambiente + manifest dati con hash;
\item
  \texttt{R3} - replay indipendente riuscito su infrastruttura esterna.
\end{enumerate}

Per claims con impatto teorico alto, il livello minimo raccomandato e \texttt{R2}; per promozione critica e preferibile \texttt{R3}.

\begin{center}\rule{0.5\linewidth}{0.5pt}\end{center}

\section{5. Proposizioni e teoremi locali}

\subsection{Proposizione 12.1 (Necessità della pre-registrazione)}

Enunciato:
Senza pre-registrazione dei criteri decisionali, la validazione OCT è esposta a bias di conferma non controllati.

Intuizione:
la decisione ex post può essere adattata al risultato osservato.

Stato: \texttt{validated}.

\subsection{Proposizione 12.2 (Insufficienza del benchmark singolo)}

Enunciato:
Un singolo benchmark non è condizione sufficiente per promuovere un claim a \texttt{validated}.

Intuizione:
la robustezza richiede indipendenza minima delle evidenze.

Stato: \texttt{validated}.

\subsection{Teorema 12.3 (Sufficienza operativa della pipeline PV-8)}

Ipotesi:

\begin{enumerate}
\def\labelenumi{\arabic{enumi}.}
\tightlist
\item
  esecuzione completa dei passi PV-8;
\item
  tracciabilità completa;
\item
  decisione conforme alla matrice standard.
\end{enumerate}

Tesi:
La decisione epistemica risultante è auditabile e non arbitraria.

Proof sketch:

\begin{enumerate}
\def\labelenumi{\arabic{enumi}.}
\tightlist
\item
  PV-8 rende esplicite tutte le dipendenze dell'esperimento;
\item
  tracciabilità completa rende verificabile il processo;
\item
  la matrice decisionale vincola la scelta finale.
\end{enumerate}

Conseguenze:
fornisce base operativa per passaggio di stato dei teoremi.

Stato: \texttt{in\_proof}.

\subsection{Teorema 12.4 (Compatibilità tra rigore formale e rigore sperimentale)}

Ipotesi:

\begin{enumerate}
\def\labelenumi{\arabic{enumi}.}
\tightlist
\item
  claim formalmente ben specificato;
\item
  protocollo PV-8 applicato;
\item
  evidenza multi-benchmark coerente.
\end{enumerate}

Tesi:
È possibile mantenere simultaneamente rigore matematico e rigore sperimentale nel ciclo OCT.

Proof sketch:

\begin{enumerate}
\def\labelenumi{\arabic{enumi}.}
\tightlist
\item
  il formalismo definisce cosa testare;
\item
  il protocollo definisce come testare;
\item
  la decisione tracciata definisce quando accettare o rivedere.
\end{enumerate}

Conseguenze:
evita la falsa alternativa tra ``teoria pura'' e ``solo empiria''.

Stato: \texttt{in\_proof}.

\subsection{Proposizione 12.5 (Monotonia informativa della tracciabilità)}

Enunciato:
A parita di esito metrico, unita con livello di tracciabilità superiore (\texttt{Trace}) produce decisioni più auditabili e meno controvertibili.

Intuizione:
la trasparenza riduce ambiguità interpretative e facilità la replica indipendente.

Stato: \texttt{validated}.

\subsection{Proposizione 12.6 (Necessità del budget d'incertezza)}

Enunciato:
Senza quantificazione esplicita del budget d'incertezza (\texttt{B\_unc}), la distinzione tra \texttt{revise} e \texttt{validated} resta metodologicamente fragile.

Intuizione:
la stessa metrica può apparire forte o debole a seconda della varianza nascosta non esplicitata.

Stato: \texttt{in\_review}.

\subsection{\texorpdfstring{Teorema 12.7 (Consistenza procedurale dell'operatore \texttt{D\_OCT})}{Teorema 12.7 (Consistenza procedurale dell'operatore D\_OCT)}}

Ipotesi:

\begin{enumerate}
\def\labelenumi{\arabic{enumi}.}
\tightlist
\item
  preregistrazione completa;
\item
  regole decisionali fissate ex ante;
\item
  applicazione uniforme su benchmark indipendenti.
\end{enumerate}

Tesi:
L'operatore \texttt{D\_OCT} e consistente nel senso procedurale: a parita di input e regole, produce decisioni replicabili tra valutatori.

Proof sketch:

\begin{enumerate}
\def\labelenumi{\arabic{enumi}.}
\tightlist
\item
  la preregistrazione blocca il perimetro decisionale;
\item
  la codifica degli input riduce ambiguità semantiche;
\item
  la matrice decisionale limita la discrezionalità ex post.
\end{enumerate}

Conseguenze:
fornisce base per audit inter-team e comparazione temporale delle decisioni.

Stato: \texttt{in\_proof}.

\begin{center}\rule{0.5\linewidth}{0.5pt}\end{center}

\section{6. Implicazioni operative}

\subsection{6.1 Per progettazione benchmark (Capitolo 13)}

Il Capitolo 13 dovrà derivare direttamente da PV-8:

\begin{enumerate}
\def\labelenumi{\arabic{enumi}.}
\tightlist
\item
  benchmark primari e secondari;
\item
  manifest dati/versioni;
\item
  criteri di conferma/falsificazione pre-registrati.
\end{enumerate}

\subsection{6.2 Per governance epistemica (Capitolo 14)}

Il Capitolo 14 userà:

\begin{enumerate}
\def\labelenumi{\arabic{enumi}.}
\tightlist
\item
  matrice decisionale standard;
\item
  regole di versionamento protocollo;
\item
  policy di pubblicazione dei fallimenti.
\end{enumerate}

\subsection{6.3 Per team di ricerca multipli}

Il protocollo consente collaborazione distribuita con linguaggio comune:

\begin{enumerate}
\def\labelenumi{\arabic{enumi}.}
\tightlist
\item
  stesso schema di report;
\item
  stessi stati decisionali;
\item
  stessa mappa di responsabilità metodologica.
\end{enumerate}

\subsection{6.4 Per pubblicazione tecnica}

Ogni release dovrebbe includere:

\begin{enumerate}
\def\labelenumi{\arabic{enumi}.}
\tightlist
\item
  versione protocollo usata;
\item
  claim valutati;
\item
  esiti e motivazioni;
\item
  limiti aperti.
\end{enumerate}

Questo migliora la leggibilità esterna e riduce ambiguità interpretative.

\subsection{6.5 Per continuità del manoscritto}

Il protocollo funziona da ``strato comune'' tra teoria (capitoli 1-11) e applicazioni/benchmark (capitoli 13-16).\\
Senza questo strato, il libro rischierebbe frammentazione in blocchi non comunicanti.

\subsection{6.6 Per pacchetto pubblicazione multi-piattaforma}

Il protocollo del Capitolo 12 consente una pubblicazione coordinata su canali diversi senza perdita di significato metodologico.

Mappatura operativa consigliata:

\begin{enumerate}
\def\labelenumi{\arabic{enumi}.}
\tightlist
\item
  GitHub: versione eseguibile (documenti, script, manifest, changelog);
\item
  Zenodo: snapshot citabile con DOI di release;
\item
  OSF: progetto di ricerca con artefatti, preregistrazioni e versioni intermedie;
\item
  Hugging Face: eventuali dataset mirror e card metodologiche per benchmark ripetibili.
\end{enumerate}

La coerenza tra piattaforme si ottiene mantenendo identico:

\begin{enumerate}
\def\labelenumi{\arabic{enumi}.}
\tightlist
\item
  identificatore di versione del protocollo;
\item
  manifest dei dataset con hash;
\item
  decision table usata per \texttt{validated/revise/reject}.
\end{enumerate}

\subsection{6.7 Per lettori non specialisti}

Un protocollo esplicito e anche una funzione pedagogica: rende leggibile il passaggio da teoria ad evidenza.

Per una fruizione ``a prova di bambino'', ogni report dovrebbe includere:

\begin{enumerate}
\def\labelenumi{\arabic{enumi}.}
\tightlist
\item
  obiettivo in linguaggio naturale;
\item
  cosa e stato misurato e perché;
\item
  esito finale e motivo della scelta;
\item
  cosa manca ancora per aumentare la confidenza.
\end{enumerate}

Questo non semplifica la teoria, ma semplifica il suo accesso.

\begin{center}\rule{0.5\linewidth}{0.5pt}\end{center}

\section{7. Limiti e non-claim}

Limiti espliciti:

\begin{enumerate}
\def\labelenumi{\arabic{enumi}.}
\tightlist
\item
  il capitolo definisce un protocollo minimo, non esaustivo per tutti i domini;
\item
  alcune metriche ausiliarie richiedono specifica per caso d'uso;
\item
  la robustezza statistica dipende dalla qualità dei dataset disponibili.
\end{enumerate}

Failure mode riconosciuti:

\begin{enumerate}
\def\labelenumi{\arabic{enumi}.}
\tightlist
\item
  pre-registrazione incompleta o retroattiva;
\item
  benchmark non indipendenti presentati come indipendenti;
\item
  decisioni non allineate ai criteri dichiarati;
\item
  omissione di fallimenti significativi nei report pubblici.
\end{enumerate}

Contromisure prioritarie:

\begin{enumerate}
\def\labelenumi{\arabic{enumi}.}
\tightlist
\item
  checklist di preregistrazione bloccante prima dell'esecuzione;
\item
  revisione incrociata minima a due valutatori per la decisione finale;
\item
  audit trimestrale dei livelli \texttt{R0-R3} sugli artefatti pubblici;
\item
  registro pubblico delle revisioni che motivi ogni cambio di stato dei claim.
\end{enumerate}

Box C - Non-claim

Questo capitolo non implica:

\begin{enumerate}
\def\labelenumi{\arabic{enumi}.}
\tightlist
\item
  che il protocollo garantisca automaticamente risultati positivi;
\item
  che ogni claim in \texttt{in\_proof} debba diventare \texttt{validated};
\item
  che la standardizzazione metodologica sostituisca l'innovazione teorica.
\end{enumerate}

\begin{center}\rule{0.5\linewidth}{0.5pt}\end{center}

\section{8. Claim Ledger (S0-S3)}

{\def\LTcaptype{none} % do not increment counter
\begin{longtable}[]{@{}lllll@{}}
\toprule\noalign{}
ID & Claim & Tipo & Evidenza & Stato \\
\midrule\noalign{}
\endhead
\bottomrule\noalign{}
\endlastfoot
C12.1 & La pipeline PV-8 è sufficiente come baseline operativa di validazione OCT & metodologico & S2 & in\_review \\
C12.2 & La pre-registrazione riduce bias decisionali in modo significativo & epistemico & S1 & validated \\
C12.3 & Multi-benchmark e multi-contesto sono necessari per promozione a \texttt{validated} & metodologico & S1 & validated \\
C12.4 & La matrice decisionale standard migliora tracciabilità e auditabilità & metodologico & S1 & in\_review \\
C12.5 & Il protocollo unifica rigore formale e rigore sperimentale & metateorico & S2 & in\_proof \\
C12.6 & La pubblicazione dei fallimenti è condizione tecnica di replicabilità & epistemico-operativo & S1 & validated \\
C12.7 & L'operatore \texttt{D\_OCT} rende replicabile la decisione a parita di input/regole & metodologico-formale & S2 & in\_proof \\
C12.8 & Il budget d'incertezza e necessario per distinguere validazione da revisione & metodologico & S1 & in\_review \\
C12.9 & I livelli \texttt{R0-R3} migliorano la comparabilità di riproducibilità tra studi & operativo & S1 & validated \\
\end{longtable}
}

\begin{center}\rule{0.5\linewidth}{0.5pt}\end{center}

\section{9. Bridge al Capitolo 13}

Il Capitolo 12 ha definito la procedura standard di validazione:

\begin{enumerate}
\def\labelenumi{\arabic{enumi}.}
\tightlist
\item
  metrica canonica e ausiliaria;
\item
  osservazione contestuale registrata;
\item
  decisione epistemica tracciata.
\end{enumerate}

Il Capitolo 13 tradurrà questa procedura in disegno benchmark completo:

\begin{enumerate}
\def\labelenumi{\arabic{enumi}.}
\tightlist
\item
  selezione e stratificazione dataset;
\item
  benchmark indipendenti e replicabilità multi-dominio;
\item
  criteri quantitativi di conferma/falsificazione per i blocchi F/D/A.
\end{enumerate}

In sintesi: il Capitolo 12 definisce il metodo; il Capitolo 13 ne implementa l'infrastruttura sperimentale.

\begin{center}\rule{0.5\linewidth}{0.5pt}\end{center}

\section{Riferimenti}

Riferimenti di framework interno:

\begin{enumerate}
\def\labelenumi{\arabic{enumi}.}
\tightlist
\item
  Ghioni, F. (2026). \texttt{TE\_CORE\_v5.1.md}.
\item
  Ghioni, F. (2026). \texttt{Ordinative\_Set\_Theory\_OST\_A\_Concise\_Guide\_For\_AI\_v2\_1.md}.
\item
  \texttt{OCT\_VALIDATION\_PROTOCOL\_v0\_1.md}.
\item
  \texttt{OCT\_FOUNDATIONAL\_BOOK\_CHAPTER\_11\_v1\_0.md}.
\item
  \texttt{OCT\_FOUNDATIONAL\_BOOK\_CHAPTER\_10\_v1\_0.md}.
\end{enumerate}

Riferimenti primari:

\begin{enumerate}
\def\labelenumi{\arabic{enumi}.}
\tightlist
\item
  Popper, K. (1959). \emph{The Logic of Scientific Discovery}.
\item
  Lakatos, I. (1978). \emph{The Methodology of Scientific Research Programmes}.
\item
  Nosek, B. et al. (2018). \emph{The preregistration revolution}.
\end{enumerate}
